\documentclass[11pt]{article}

\usepackage{xcolor}
\usepackage{scrlayer-scrpage}

\newcommand{\MyName}{P. Lerdkijrachapong}

\usepackage{blindtext}

\clearscrheadfoot
\rohead{\MyName}
\lehead{\today}

\chead{Motivation}
\lohead{07 September, 2025}
\rehead{\MyName}    
\cefoot{\thepage}

\begin{document}

Before we get to implement, use, or weaponise ``motivation", we first have to understand what motivation is, and what are theories related to it.

The word Motivation derived from Latin ``\textit{Movere}", meaning ``\textit{to move}". As the meaning evolve, motivation is the willingness to exert high levels of effort toward a goal. And as there're concepts, there're theories related to it.

The main theories of motivation can be categorised in to two part: 1) Content Theories, and 2) Process Theories. To quickly understand the core concept of each theory type: The Content Theories identify the \textbf{needs} that motivate people's behaviour, while The Process Theories mainly focus on the \textbf{thought processes} that act to motivate people's behaviour.

The widely spread Content Theories are:
\begin{enumerate}
	\item Maslow's Hierarchy of Needs:

Basic physiological needs to self-actualisation and transcendence.

	\item Herzberg's Two-Factor Theory:
		\begin{itemize}
			\item Hygiene Factor (e.g., salary, work conditions) \textbf{prevent dissatisfaction}.
			\item Motivators (e.g., achievement, recognition) \textbf{drive satisfaction}.
		\end{itemize}
	\item McClelland's Acquired Needs Theory:
		\begin{itemize}
			\item The Need for Achievement, Affiliation, and Power.
		\end{itemize}
	\item McGregor's Theory X and Theory Y:
		\begin{itemize}
			\item Theory X: People dislike work and need control.
			\item Theory Y: People are self-motivated and seek responsibility.
		\end{itemize}
\end{enumerate}
and the widely known Process Theories are:
\begin{enumerate}
	\item Expectancy Theory by \textit{V. Vroom}:
	\begin{center}
			Motivation $\equiv Expectancy\times Instrumentality \times {Valence} .$
	\end{center}

	\item Equity Theory by \textit{J.S. Adams}:

		Motivation is influenced by perceived fairness in input\textendash output ratios.
	\item Goal Setting Theory by \textit{Edwin A. Locke \& Gary P. Latham}:

Specific, challenging goals with feedback enhance performance.
	\item Reinforcement Theory by \textit{B.F. Skinner}:
	
Behavior is shaped by rewards and punishments.
\end{enumerate}


By understanding Motivation Theories, we can better understand what drives employees from different cultures, and improve Cross-Cultural Management within the organisation. Moreover, we can apply Motivation Theories to enhance employees engagement by understood their value, and opinion. For example:

Managing a Diverse Team Project:
\begin{itemize}
	\item By using \textbf{Expectancy Theory}.
	\begin{itemize}
		\item Set Clear Goal (\textbf{Expectancy})
		\item Set Transparent Reward System (\textbf{Instrumentality})
		\item Reward the member (\textbf{Valence})
	\end{itemize}
\end{itemize}
As you can see from the demonstration above, by implementing Motivation Theory, we can improve the workplace management, and enhance employee's engagement,
\end{document}
